\section{Discussion}
\label{sec:discussion}

\subsection{Why Confirmation Resists Learning}

The two negative results of \cref{sec:inversion} and \cref{sec:learned} have a single
explanation, and it is a property of the data rather than of the methods. The confirmation
stage receives roughly two hundred positive tracks against 7000--28,000 negatives, and the
negatives are not a homogeneous class. They divide into false candidates, which are easy,
and \emph{duplicate} tracks on animals already counted, which are not: a duplicate is a
correct detection of a real animal, and \cref{sec:analysis} shows that duplicates sit
between the positives and the false candidates on every feature we measured: box size, track length, candidate confidence alike.

A confirmer therefore receives contradictory supervision on exactly the examples that
distinguish counting from detection. That is why capacity does not help: the capacity trend in \cref{sec:learned} declines monotonically from 13 parameters to $6\times10^4$, and
shrinking a gradient-boosted model to depth-1 stumps recovers nearly the whole gap. It is also why every improvement to candidate generation costs count. Growing the pool grows the
ambiguous middle of the negative class far faster than it grows the positives, so a rule
tuned to survive the new ratio must be more aggressive, and aggression costs real animals.
Extrapolating the capacity trend puts the crossover in the low thousands of positive
tracks, an order of magnitude beyond this corpus, and a concrete target for the next
dataset rather than for a better architecture.

\subsection{Recall and Duplicate Suppression Are Bought with One Parameter}

\Cref{sec:sensitivity} makes the same mechanism visible in a single parameter and is, we
think, the most transferable observation in the paper. Deleting the confidence condition
raises held-out coverage by 23 points, more than any other change we made, and simultaneously
turns a systematic under-count into a systematic over-count, because roughly two duplicates
are admitted for each genuine animal recovered. The condition is not filtering non-animals.
It is the only cheap signal separating a first sighting from a re-sighting, and the recall it
appears to be discarding is the price of that separation.

This reframes what a practitioner should tune. The question is not ``how low can the
confidence threshold go before false positives appear'', the usual detection framing, under
which $c=0$ looks free because coverage rises and the extra tracks are all on real animals. It
is ``at what point does the estimator's bias cross zero'', a different quantity with a
different optimum, which the next subsection locates and costs.

\subsection{Choosing an Operating Point for Abundance Estimation}
\label{sec:operating}

A survey ecologist and a per-transect report want different things from this system, and the
confirmation threshold is where the difference is settled.

Our headline rule minimises mean absolute error per video, and it under-counts: 58 animals
predicted where 83 exist across the 13 held-out transects, a shortfall of 25. For an abundance estimate that is not the quantity to optimise. Setting $c=0.45$ instead, at the point where signed bias crosses zero in \cref{tab:confsens}, makes the same pipeline predict \textbf{81 animals against 83}, an aggregate error of 2.4\%. Individual transects are
less accurate, with mean absolute error rising from 2.38 to 3.38 animals per video, but the
errors are now two-sided and cancel across a survey rather than accumulating in one direction.
Under-counting by 25 becomes over- and under-counting that nets to two.

Which is preferable depends on the downstream statistic, and the two are not
interchangeable. A management unit's population total is a sum over transects, so unbiasedness
matters and per-transect variance largely averages out; a report of how many animals were seen
on a particular road on a particular night is a per-transect quantity, and there the
low-variance rule is correct. We report the minimum-error rule as our headline because it is
the conventional choice for a counting system, and record the unbiased setting here because it
is the one an abundance study should use.

This distinction is obscured by our own headline metric. The capped fraction of
\cref{sec:decomposition} deliberately prevents an over-count in one video from offsetting a
miss in another, which suits a per-video counter but not a population total, where that cancellation is the point. Under the capped metric $c=0.45$ scores
72.3\%; under the aggregate it recovers 97.6\% of the animals present.

One route we expected to work does not. The standard remedy for a known undercount is a
sightability correction: estimate the detected fraction and divide by it. The held-out
protocol makes that estimate impossible to obtain honestly from our own data. On the 19
detector-training videos the pipeline predicts 152 animals against 153, so a correction factor
fitted there is 1.007, effectively no correction, while the held-out videos need 1.43. The
contamination that inflates coverage on seen video also destroys any correction estimated from
it. A sightability factor for this system would have to be fitted on video the detector never
trained on, which is to say on data that could have been used to train it instead.

A second limit on any abundance estimate lies upstream of the pipeline, in the ground truth.
As in any visual survey, an annotated animal is one a reviewer could see in the footage, so
every figure in this paper measures the pipeline against what is visible rather than against
the animals present. We tested how strongly this depends on range, with box size as the range
proxy of \cref{sec:analysis}. If deer were spread evenly over the ground in view, the number of
animals per size bin would grow steeply as boxes shrink, because the area in view grows with
distance. The 235 annotated animals follow that law closely between 30 and 60\,px but fall
short of it below about 25\,px: 29 animals are annotated at 15--20\,px where the geometry
predicts about 210, and about 90 even under the most conservative assumption, a fixed-width
strip along the road. The annotated tracks do not shorten significantly with range
(Kruskal--Wallis $p=0.24$ across size quartiles), so a distant animal, once visible, is
followed about as long as a near one. The likely cause is what the sensor can resolve: at
$640\times512$ a distant deer is a few warm pixels, often partly behind vegetation or terrain,
the same fall in detectability with range that spotlight counts
face~\cite{vercauteren2011managing}. The comparison assumes deer are spread evenly with
distance from the road, which a road transect cannot guarantee~\cite{diefenbach2025accounting},
and box size also varies with body size and posture. An abundance estimate built on this
system would therefore need both factors: the pipeline's size curve of \cref{sec:analysis}
and a detection function for the survey itself.

\subsection{Implications for Practice}

What follows are consequences of the measurements above rather than recommendations drawn
from survey experience, and they are scoped to what this corpus can support: one sensor, four
sites, one winter, one species, 236 individuals. Within that scope, four of them run against
the field's default.

Selecting a detector by its detection metric did not select the model that counts best here.
The counting criterion re-ranks seven of eleven models by two or more places
(\cref{sec:choosing}), and the best model by AP$_{50}$ counted worse than the fourth-best.
Evaluating on video the detector had trained on inflated coverage by 15.5 points
(\cref{tab:leak}), which is what makes a held-out protocol worth its cost. A larger candidate
pool was not a better one in any of the four attempts we made (\cref{tab:ablationmain}). And a learned confirmation stage did not beat a three-parameter rule at this data scale, with the capacity
trend placing the crossover an order of magnitude beyond 236 individuals
(\cref{tab:confirmers}).

Three observations point the other way. Contrast normalization was worth more than any
architectural choice we tested, 0.519 against 0.299 test mAP50, a larger effect than the entire 0.365--0.510 spread across eleven architectures. The upper tail of the object-scale
distribution repaid the attention normally reserved for the lower one: three of our misses are
animals \emph{larger} than anything in the training split (\cref{fig:scale}), a failure mode a
small-object-focused evaluation would not surface. And an absolute confidence threshold did not
transfer between sites, since a detector trained without a site scored that site systematically lower and a threshold
fitted elsewhere rejected every track on one fold (\cref{sec:loso}), a shift this corpus cannot
separate from the different nights on which the sites were surveyed.
Whether a small labelled sample from each new site or survey night is the right remedy is a question this
corpus can pose but not answer.

\subsection{Limitations}
\label{sec:limitations}

Several limits bound these claims. The ground truth contains only animals visible in the
footage, and it thins out with range (\cref{sec:operating}), so every figure measures the
pipeline against what the video shows rather than against the animals present. The corpus is
small and narrow: 236 animals of one species from four sites in one winter, each site
surveyed on a single night so that site and night are confounded, one site contributing
56\% of the animals, and a single tracker (BoT-SORT). Four of the 13 held-out videos are the
detector's validation split and influenced checkpoint selection (\cref{sec:heldout}). The
headline count is a per-video aggregate; the identity-matched figure is 10 points lower. No
animal appears in more than one video, so cross-session identification could not be
evaluated. Finally, the confidence threshold does not transfer across sites
(\cref{sec:loso}), so deployment elsewhere requires re-fitting it on labelled local data.
